%!TeX root=../persuasiontop.tex
\chapter[Chapter \thechapter]{}

\begin{figure}[t!]
\centering
\includegraphics[width=\linewidth]{nouveau10}
\end{figure}

\lettrine[lines=4,lraise=0.3]{O}{ther} opportunities of making her observations could not fail to occur. Anne had soon been in company with all the four together often enough to have an opinion, though too wise to acknowledge as much at home, where she knew it would have satisfied neither husband nor wife; for while she considered Louisa to be rather the favourite, she could not but think, as far as she might dare to judge from memory and experience, that Captain Wentworth was not in love with either. They were more in love with him; yet there it was not love. It was a little fever of admiration; but it might, probably must, end in love with some. Charles Hayter seemed aware of being slighted, and yet Henrietta had sometimes the air of being divided between them. Anne longed for the power of representing to them all what they were about, and of pointing out some of the evils they were exposing themselves to. She did not attribute guile to any. It was the highest satisfaction to her to believe Captain Wentworth not in the least aware of the pain he was occasioning. There was no triumph, no pitiful triumph in his manner. He had, probably, never heard, and never thought of any claims of Charles Hayter. He was only wrong in accepting the attentions (for accepting must be the word) of two young women at once.

After a short struggle, however, Charles Hayter seemed to quit the field. Three days had passed without his coming once to Uppercross; a most decided change. He had even refused one regular invitation to dinner; and having been found on the occasion by Mr Musgrove with some large books before him, Mr and Mrs Musgrove were sure all could not be right, and talked, with grave faces, of his studying himself to death. It was Mary's hope and belief that he had received a positive dismissal from Henrietta, and her husband lived under the constant dependence of seeing him to-morrow. Anne could only feel that Charles Hayter was wise.

One morning, about this time Charles Musgrove and Captain Wentworth being gone a-shooting together, as the sisters in the Cottage were sitting quietly at work, they were visited at the window by the sisters from the Mansion-house.

It was a very fine November day, and the Miss Musgroves came through the little grounds, and stopped for no other purpose than to say, that they were going to take a \textit{long} walk, and, therefore, concluded Mary could not like to go with them; and when Mary immediately replied, with some jealousy at not being supposed a good walker, »Oh, yes, I should like to join you very much, I am very fond of a long walk;« Anne felt persuaded, by the looks of the two girls, that it was precisely what they did not wish, and admired again the sort of necessity which the family habits seemed to produce, of everything being to be communicated, and everything being to be done together, however undesired and inconvenient. She tried to dissuade Mary from going, but in vain; and that being the case, thought it best to accept the Miss Musgroves' much more cordial invitation to herself to go likewise, as she might be useful in turning back with her sister, and lessening the interference in any plan of their own.

»I cannot imagine why they should suppose I should not like a long walk,« said Mary, as she went up stairs. »Everybody is always supposing that I am not a good walker; and yet they would not have been pleased, if we had refused to join them. When people come in this manner on purpose to ask us, how can one say no?«

Just as they were setting off, the gentlemen returned. They had taken out a young dog, who had spoilt their sport, and sent them back early. Their time and strength, and spirits, were, therefore, exactly ready for this walk, and they entered into it with pleasure. Could Anne have foreseen such a junction, she would have staid at home; but, from some feelings of interest and curiosity, she fancied now that it was too late to retract, and the whole six set forward together in the direction chosen by the Miss Musgroves, who evidently considered the walk as under their guidance.

Anne's object was, not to be in the way of anybody; and where the narrow paths across the fields made many separations necessary, to keep with her brother and sister. Her \textit{pleasure} in the walk must arise from the exercise and the day, from the view of the last smiles of the year upon the tawny leaves, and withered hedges, and from repeating to herself some few of the thousand poetical descriptions extant of autumn, that season of peculiar and inexhaustible influence on the mind of taste and tenderness, that season which had drawn from every poet, worthy of being read, some attempt at description, or some lines of feeling. She occupied her mind as much as possible in such like musings and quotations; but it was not possible, that when within reach of Captain Wentworth's conversation with either of the Miss Musgroves, she should not try to hear it; yet she caught little very remarkable. It was mere lively chat, such as any young persons, on an intimate footing, might fall into. He was more engaged with Louisa than with Henrietta. Louisa certainly put more forward for his notice than her sister. This distinction appeared to increase, and there was one speech of Louisa's which struck her. After one of the many praises of the day, which were continually bursting forth, Captain Wentworth added:—

»What glorious weather for the Admiral and my sister! They meant to take a long drive this morning; perhaps we may hail them from some of these hills. They talked of coming into this side of the country. I wonder whereabouts they will upset to-day. Oh! it does happen very often, I assure you; but my sister makes nothing of it; she would as lieve be tossed out as not.«

»Ah! You make the most of it, I know,« cried Louisa, »but if it were really so, I should do just the same in her place. If I loved a man, as she loves the Admiral, I would always be with him, nothing should ever separate us, and I would rather be overturned by him, than driven safely by anybody else.«

It was spoken with enthusiasm.

»Had you?« cried he, catching the same tone; »I honour you!« And there was silence between them for a little while.

Anne could not immediately fall into a quotation again. The sweet scenes of autumn were for a while put by, unless some tender sonnet, fraught with the apt analogy of the declining year, with declining happiness, and the images of youth and hope, and spring, all gone together, blessed her memory. She roused herself to say, as they struck by order into another path, »Is not this one of the ways to Winthrop?« But nobody heard, or, at least, nobody answered her.

Winthrop, however, or its environs—for young men are, sometimes to be met with, strolling about near home—was their destination; and after another half mile of gradual ascent through large enclosures, where the ploughs at work, and the fresh made path spoke the farmer counteracting the sweets of poetical despondence, and meaning to have spring again, they gained the summit of the most considerable hill, which parted Uppercross and Winthrop, and soon commanded a full view of the latter, at the foot of the hill on the other side.

Winthrop, without beauty and without dignity, was stretched before them; an indifferent house, standing low, and hemmed in by the barns and buildings of a farm-yard.

Mary exclaimed, »Bless me! here is Winthrop. I declare I had no idea! Well now, I think we had better turn back; I am excessively tired.«

Henrietta, conscious and ashamed, and seeing no cousin Charles walking along any path, or leaning against any gate, was ready to do as Mary wished; but »No!« said Charles Musgrove, and »No, no!« cried Louisa more eagerly, and taking her sister aside, seemed to be arguing the matter warmly.

Charles, in the meanwhile, was very decidedly declaring his resolution of calling on his aunt, now that he was so near; and very evidently, though more fearfully, trying to induce his wife to go too. But this was one of the points on which the lady shewed her strength; and when he recommended the advantage of resting herself a quarter of an hour at Winthrop, as she felt so tired, she resolutely answered, »Oh! no, indeed! walking up that hill again would do her more harm than any sitting down could do her good;« and, in short, her look and manner declared, that go she would not.

After a little succession of these sort of debates and consultations, it was settled between Charles and his two sisters, that he and Henrietta should just run down for a few minutes, to see their aunt and cousins, while the rest of the party waited for them at the top of the hill. Louisa seemed the principal arranger of the plan; and, as she went a little way with them, down the hill, still talking to Henrietta, Mary took the opportunity of looking scornfully around her, and saying to Captain Wentworth—

»It is very unpleasant, having such connexions! But, I assure you, I have never been in the house above twice in my life.«

She received no other answer, than an artificial, assenting smile, followed by a contemptuous glance, as he turned away, which Anne perfectly knew the meaning of.

The brow of the hill, where they remained, was a cheerful spot: Louisa returned; and Mary, finding a comfortable seat for herself on the step of a stile, was very well satisfied so long as the others all stood about her; but when Louisa drew Captain Wentworth away, to try for a gleaning of nuts in an adjoining hedge-row, and they were gone by degrees quite out of sight and sound, Mary was happy no longer; she quarrelled with her own seat, was sure Louisa had got a much better somewhere, and nothing could prevent her from going to look for a better also. She turned through the same gate, but could not see them. Anne found a nice seat for her, on a dry sunny bank, under the hedge-row, in which she had no doubt of their still being, in some spot or other. Mary sat down for a moment, but it would not do; she was sure Louisa had found a better seat somewhere else, and she would go on till she overtook her.

Anne, really tired herself, was glad to sit down; and she very soon heard Captain Wentworth and Louisa in the hedge-row, behind her, as if making their way back along the rough, wild sort of channel, down the centre. They were speaking as they drew near. Louisa's voice was the first distinguished. She seemed to be in the middle of some eager speech. What Anne first heard was—

»And so, I made her go. I could not bear that she should be frightened from the visit by such nonsense. What! would I be turned back from doing a thing that I had determined to do, and that I knew to be right, by the airs and interference of such a person, or of any person I may say? No, I have no idea of being so easily persuaded. When I have made up my mind, I have made it; and Henrietta seemed entirely to have made up hers to call at Winthrop to-day; and yet, she was as near giving it up, out of nonsensical complaisance!«

»She would have turned back then, but for you?«

»She would indeed. I am almost ashamed to say it.«

»Happy for her, to have such a mind as yours at hand! After the hints you gave just now, which did but confirm my own observations, the last time I was in company with him, I need not affect to have no comprehension of what is going on. I see that more than a mere dutiful morning visit to your aunt was in question; and woe betide him, and her too, when it comes to things of consequence, when they are placed in circumstances requiring fortitude and strength of mind, if she have not resolution enough to resist idle interference in such a trifle as this. Your sister is an amiable creature; but \textit{yours} is the character of decision and firmness, I see. If you value her conduct or happiness, infuse as much of your own spirit into her as you can. But this, no doubt, you have been always doing. It is the worst evil of too yielding and indecisive a character, that no influence over it can be depended on. You are never sure of a good impression being durable; everybody may sway it. Let those who would be happy be firm. Here is a nut,« said he, catching one down from an upper bough, »to exemplify: a beautiful glossy nut, which, blessed with original strength, has outlived all the storms of autumn. Not a puncture, not a weak spot anywhere. This nut,« he continued, with playful solemnity, »while so many of his brethren have fallen and been trodden under foot, is still in possession of all the happiness that a hazel nut can be supposed capable of.« Then returning to his former earnest tone—»My first wish for all whom I am interested in, is that they should be firm. If Louisa Musgrove would be beautiful and happy in her November of life, she will cherish all her present powers of mind.«

He had done, and was unanswered. It would have surprised Anne if Louisa could have readily answered such a speech: words of such interest, spoken with such serious warmth! She could imagine what Louisa was feeling. For herself, she feared to move, lest she should be seen. While she remained, a bush of low rambling holly protected her, and they were moving on. Before they were beyond her hearing, however, Louisa spoke again.

»Mary is good-natured enough in many respects,« said she; »but she does sometimes provoke me excessively, by her nonsense and pride—the Elliot pride. She has a great deal too much of the Elliot pride. We do so wish that Charles had married Anne instead. I suppose you know he wanted to marry Anne?«

After a moment's pause, Captain Wentworth said—

»Do you mean that she refused him?«

»Oh! yes; certainly.«

»When did that happen?«

»I do not exactly know, for Henrietta and I were at school at the time; but I believe about a year before he married Mary. I wish she had accepted him. We should all have liked her a great deal better; and papa and mamma always think it was her great friend Lady Russell's doing, that she did not. They think Charles might not be learned and bookish enough to please Lady Russell, and that therefore, she persuaded Anne to refuse him.«

The sounds were retreating, and Anne distinguished no more. Her own emotions still kept her fixed. She had much to recover from, before she could move. The listener's proverbial fate was not absolutely hers; she had heard no evil of herself, but she had heard a great deal of very painful import. She saw how her own character was considered by Captain Wentworth, and there had been just that degree of feeling and curiosity about her in his manner which must give her extreme agitation.

As soon as she could, she went after Mary, and having found, and walked back with her to their former station, by the stile, felt some comfort in their whole party being immediately afterwards collected, and once more in motion together. Her spirits wanted the solitude and silence which only numbers could give.

Charles and Henrietta returned, bringing, as may be conjectured, Charles Hayter with them. The minutiae of the business Anne could not attempt to understand; even Captain Wentworth did not seem admitted to perfect confidence here; but that there had been a withdrawing on the gentleman's side, and a relenting on the lady's, and that they were now very glad to be together again, did not admit a doubt. Henrietta looked a little ashamed, but very well pleased;—Charles Hayter exceedingly happy: and they were devoted to each other almost from the first instant of their all setting forward for Uppercross.

Everything now marked out Louisa for Captain Wentworth; nothing could be plainer; and where many divisions were necessary, or even where they were not, they walked side by side nearly as much as the other two. In a long strip of meadow land, where there was ample space for all, they were thus divided, forming three distinct parties; and to that party of the three which boasted least animation, and least complaisance, Anne necessarily belonged. She joined Charles and Mary, and was tired enough to be very glad of Charles's other arm; but Charles, though in very good humour with her, was out of temper with his wife. Mary had shewn herself disobliging to him, and was now to reap the consequence, which consequence was his dropping her arm almost every moment to cut off the heads of some nettles in the hedge with his switch; and when Mary began to complain of it, and lament her being ill-used, according to custom, in being on the hedge side, while Anne was never incommoded on the other, he dropped the arms of both to hunt after a weasel which he had a momentary glance of, and they could hardly get him along at all.

This long meadow bordered a lane, which their footpath, at the end of it was to cross, and when the party had all reached the gate of exit, the carriage advancing in the same direction, which had been some time heard, was just coming up, and proved to be Admiral Croft's gig. He and his wife had taken their intended drive, and were returning home. Upon hearing how long a walk the young people had engaged in, they kindly offered a seat to any lady who might be particularly tired; it would save her a full mile, and they were going through Uppercross. The invitation was general, and generally declined. The Miss Musgroves were not at all tired, and Mary was either offended, by not being asked before any of the others, or what Louisa called the Elliot pride could not endure to make a third in a one horse chaise.

The walking party had crossed the lane, and were surmounting an opposite stile, and the Admiral was putting his horse in motion again, when Captain Wentworth cleared the hedge in a moment to say something to his sister. The something might be guessed by its effects.

»Miss Elliot, I am sure \textit{you} are tired,« cried Mrs Croft. »Do let us have the pleasure of taking you home. Here is excellent room for three, I assure you. If we were all like you, I believe we might sit four. You must, indeed, you must.«

Anne was still in the lane; and though instinctively beginning to decline, she was not allowed to proceed. The Admiral's kind urgency came in support of his wife's; they would not be refused; they compressed themselves into the smallest possible space to leave her a corner, and Captain Wentworth, without saying a word, turned to her, and quietly obliged her to be assisted into the carriage.

Yes; he had done it. She was in the carriage, and felt that he had placed her there, that his will and his hands had done it, that she owed it to his perception of her fatigue, and his resolution to give her rest. She was very much affected by the view of his disposition towards her, which all these things made apparent. This little circumstance seemed the completion of all that had gone before. She understood him. He could not forgive her, but he could not be unfeeling. Though condemning her for the past, and considering it with high and unjust resentment, though perfectly careless of her, and though becoming attached to another, still he could not see her suffer, without the desire of giving her relief. It was a remainder of former sentiment; it was an impulse of pure, though unacknowledged friendship; it was a proof of his own warm and amiable heart, which she could not contemplate without emotions so compounded of pleasure and pain, that she knew not which prevailed.

Her answers to the kindness and the remarks of her companions were at first unconsciously given. They had travelled half their way along the rough lane, before she was quite awake to what they said. She then found them talking of »Frederick.«

»He certainly means to have one or other of those two girls, Sophy,« said the Admiral; »but there is no saying which. He has been running after them, too, long enough, one would think, to make up his mind. Ay, this comes of the peace. If it were war now, he would have settled it long ago. We sailors, Miss Elliot, cannot afford to make long courtships in time of war. How many days was it, my dear, between the first time of my seeing you and our sitting down together in our lodgings at North Yarmouth?«

»We had better not talk about it, my dear,« replied Mrs Croft, pleasantly; »for if Miss Elliot were to hear how soon we came to an understanding, she would never be persuaded that we could be happy together. I had known you by character, however, long before.«

»Well, and I had heard of you as a very pretty girl, and what were we to wait for besides? I do not like having such things so long in hand. I wish Frederick would spread a little more canvass, and bring us home one of these young ladies to Kellynch. Then there would always be company for them. And very nice young ladies they both are; I hardly know one from the other.«

»Very good humoured, unaffected girls, indeed,« said Mrs Croft, in a tone of calmer praise, such as made Anne suspect that her keener powers might not consider either of them as quite worthy of her brother; »and a very respectable family. One could not be connected with better people. My dear Admiral, that post! we shall certainly take that post.«

But by coolly giving the reins a better direction herself they happily passed the danger; and by once afterwards judiciously putting out her hand they neither fell into a rut, nor ran foul of a dung-cart; and Anne, with some amusement at their style of driving, which she imagined no bad representation of the general guidance of their affairs, found herself safely deposited by them at the Cottage.